
        \documentclass[fleqn]{article}      
        \usepackage{amsmath}
        \usepackage{fontspec}
        \usepackage{kotex} % 한국어 지원  

        \begin{document}      
        \textbf{쉬운 문제}

\noindent 1. 다음 행렬의 합을 구하세요. 
\[
A = \begin{pmatrix}
1 & 2 \\
3 & 4
\end{pmatrix}, \quad B = \begin{pmatrix}
5 & 6 \\
7 & 8
\end{pmatrix}
\] 
\ \\\\

\noindent 2. 행렬 \(A\)의 곱셈을 구하세요.
\[
A = \begin{pmatrix}
2 & 3 \\
4 & 5
\end{pmatrix}, \quad B = \begin{pmatrix}
1 & 0 \\
0 & 1
\end{pmatrix}
\]
\ \\\\

\noindent 3. 행렬 \(A\)의 전치 행렬을 구하세요.
\[
A = \begin{pmatrix}
1 & 2 & 3 \\
4 & 5 & 6
\end{pmatrix}
\]
\ \\\\

\noindent 4. 다음 행렬의 결정식을 구하세요.
\[
A = \begin{pmatrix}
2 & 3 \\
1 & 4
\end{pmatrix}
\]
\ \\\\

\noindent 5. 다음 행렬의 크기를 구하세요.
\[
A = \begin{pmatrix}
1 & 2 & 3 \\
4 & 5 & 6
\end{pmatrix}
\]
\ \\\\

\textbf{중간 문제}

\noindent 6. \(2 \times 2\) 행렬 \(A\)와 \(B\)가 주어졌을 때, 다음을 계산하세요.
\[
A = \begin{pmatrix}
1 & 2 \\
3 & 4
\end{pmatrix}, \quad B = \begin{pmatrix}
5 & 6 \\
7 & 8
\end{pmatrix}
\]
행렬 \(C = A + B\)와 \(D = A \cdot B\) 
\ \\\\

\noindent 7. 다음의 행렬 연립 방정식을 풀어라.
\[
A = \begin{pmatrix}
1 & 1 \\
2 & 3
\end{pmatrix}, \quad \begin{pmatrix}
x \\
y
\end{pmatrix} = \begin{pmatrix}
5 \\
11
\end{pmatrix}
\]
\ \\\\

\noindent 8. 다음 두 행렬의 차를 계산하세요.
\[
A = \begin{pmatrix}
3 & 2 \\
1 & 4
\end{pmatrix}, \quad B = \begin{pmatrix}
1 & 0 \\
2 & 3
\end{pmatrix}
\]
\ \\\\

\noindent 9. \(A\)와 \(B\) 행렬이 다음과 같을 때, \(B\)의 해를 구하세요.
\[
A = \begin{pmatrix}
1 & 2 \\
3 & 4
\end{pmatrix}, \quad \text{det}(A) = 1, \quad B = \begin{pmatrix}
5 \\
6
\end{pmatrix}
\]
\ \\\\

\noindent 10. \(2 \times 2\) 행렬 \(A\)의 역행렬을 구하세요. 
\[
A = \begin{pmatrix}
2 & 3 \\
5 & 7
\end{pmatrix}
\]
\ \\\\

\textbf{어려운 문제}

\noindent 11. 행렬 \(A\)의 고유값을 구하세요.
\[
A = \begin{pmatrix}
4 & 2 \\
1 & 3
\end{pmatrix}
\]
\ \\\\

\noindent 12. 아래 행렬의 행렬식과 역행렬을 구하세요.
\[
A = \begin{pmatrix}
3 & 5 \\
2 & 4
\end{pmatrix}
\]
\ \\\\

\noindent 13. 다음의 선형 방정식을 행렬을 이용하여 풀어라.
\[
\begin{pmatrix}
2 & 1 \\
1 & -1
\end{pmatrix} \begin{pmatrix}
x \\
y
\end{pmatrix} = \begin{pmatrix}
3 \\
1
\end{pmatrix}
\]
\ \\\\

\noindent 14. 두 행렬의 곱이 주어질 때, 그 크기를 구하세요.
\[
A = \begin{pmatrix}
1 & 2 \\
3 & 4 \\
5 & 6
\end{pmatrix}, \quad B = \begin{pmatrix}
7 & 8 \\
9 & 10
\end{pmatrix}
\]
\ \\\\

\noindent 15. 다음의 고유벡터를 구하세요.
\[
A = \begin{pmatrix}
2 & 1 \\
1 & 2
\end{pmatrix}
\]
\ \\\\

\textbf{답안}

1. \(\begin{pmatrix} 6 & 8 \\ 10 & 12 \end{pmatrix}\) \\\\
2. \(\begin{pmatrix} 2 & 3 \\ 4 & 5 \end{pmatrix}\) \\\\
3. \(\begin{pmatrix} 1 & 4 \\ 2 & 5 \\ 3 & 6 \end{pmatrix}\) \\\\
4. \(5\) \\\\
5. \(2 \times 3\) \\\\
6. \(C = \begin{pmatrix} 6 & 8 \\ 10 & 12 \end{pmatrix}, D = \begin{pmatrix} 19 & 22 \\ 43 & 50 \end{pmatrix}\) \\\\
7. \(x = 4, y = 1\) \\\\
8. \(A - B = \begin{pmatrix} 2 & 2 \\ -1 & 1 \end{pmatrix}\) \\\\
9. \(x = 3, y = 0\) \\\\
10. \(A^{-1} = \begin{pmatrix} 7 & -3 \\ -5 & 2 \end{pmatrix}\) \\\\
11. 고유값: \(5, 2\) \\\\
12. 행렬식: \(2, A^{-1} = \begin{pmatrix} 2 & -2.5 \\ 1 & 1.5 \end{pmatrix}\) \\\\
13. \(x = 1, y = 1\) \\\\
14. \(3 \times 2\) \\\\
15. 고유벡터: \(\begin{pmatrix} 1 \\ -1 \end{pmatrix}, \begin{pmatrix} 1 \\ 1 \end{pmatrix}\) \\\\      
        \end{document} 
    